\chapter{Results and Discussion}

% ====================================================================
\section{Overview}
% ====================================================================

This chapter presents the evaluation of the HealthFraudMLChain
pipeline, covering every stage from raw data ingestion and blockchain
registration through hyperparameter optimisation, model selection, statistical
validation, and post-hoc explainability. All numerical results reported here
are drawn directly from controlled experiments on the CMS Medicare Part-B
dataset (190 engineered features, 5{,}410 provider records). Section 4.2
discusses each sub-system in turn: the Flask dashboard and blockchain layer,
Optuna-driven tuning, 10-fold cross-validation metrics, statistical significance
tests, threshold calibration, bootstrap confidence intervals, SMOTE-ENN
ablation, feature importance, SHAP waterfall case studies, LIME local
explanations, feature-selection ablation, learning curves, and calibration
analysis. A key finding is that the weighted ensemble of gradient-boosted
models, evaluated at a fixed decision threshold of $t = 0.444$, achieves
the best overall generalisation, with a bootstrapped PR-AUC of
$0.8115\;[0.7801,\,0.8406]$ and MCC of $0.7148\;[0.6818,\,0.7474]$.

% ====================================================================
\section{Results Analysis}
% ====================================================================

% --------------------------------------------------------------------
\subsection{Data Input and Blockchain Results}
% --------------------------------------------------------------------

The system exposes a Flask-based web dashboard that serves two primary
interfaces: a \emph{prediction interface} for submitting individual or
batch provider records, and an \emph{analytics dashboard} displaying
aggregate statistics, risk-level breakdowns, and chain-health indicators.
On submission, a provider record is passed through the trained weighted
ensemble; the returned \texttt{fraud\_probability} and \texttt{risk\_level}
(\textit{Low}/\textit{Medium}/\textit{High}) are then committed to an
append-only blockchain.

Each block in the chain stores the provider identifier under ECIES
(Elliptic-Curve Integrated Encryption Scheme) public-key encryption,
ensuring that personally identifiable information is never stored in
plaintext. The chain uses SHA-256 hashing with Proof-of-Work at
difficulty level 2, and every block records its own \texttt{hash} and
the previous block's hash to maintain tamper-evidence. A Merkle tree is
constructed over each block's transaction set, with the root digest
stored as \texttt{merkle\_root}; membership proofs are verified in
$O(\log n)$ time across 6 levels.

Listing~\ref{lst:block} shows a representative JSON block from the
persisted chain, illustrating the full field layout.

\vspace{0.5em}
\noindent
\begin{minipage}{\linewidth}
\begin{lstlisting}[
  caption={Sample blockchain block - provider prediction record},
  label={lst:block},
  basicstyle=\ttfamily\footnotesize,
  breaklines=true,
  frame=single,
  numbers=left,
  numberstyle=\tiny,
  aboveskip=0.5em,
  belowskip=0.5em
]
{
  "index": 42,
  "timestamp": "2024-12-01T10:23:45.123456",
  "provider_id": "BIIyNkF3...encryptedECIES...==",
  "fraud_probability": 0.7843,
  "prediction": 1,
  "risk_level": "High",
  "previous_hash": "0082a3f1e9c4b...sha256...",
  "nonce": 178,
  "merkle_root": "a1b2c3d4e5f6...sha256...",
  "hash": "009fe2cc71bd...sha256...",
  "transactions": [{
    "provider_id_encrypted": "BIIyNkF3...==",
    "fraud_probability": 0.7843,
    "prediction": 1,
    "risk_level": "High",
    "timestamp": "2024-12-01T10:23:45.123456"
  }]
}
\end{lstlisting}
\end{minipage}

The final chain comprised 110 blocks covering all 5{,}410 provider
records. Full verification results are detailed in
Section~\ref{sec:blockchain_verification}.

% --------------------------------------------------------------------
\subsection{Optuna Hyperparameter Optimisation}
% --------------------------------------------------------------------

Hyperparameters were tuned using the Optuna framework with 60 trials per
model (except LightGBM, which triggered early pruning and completed only
1 effective trial), a 5-fold inner cross-validation loop, and
Area Under the Precision-Recall Curve (AUC-PR) as the objective. No
SMOTE oversampling was applied during the tuning phase. The complete
tuning summary is reported in Table~\ref{tab:optuna}.

\begin{table}[H]
\caption{Optuna hyperparameter optimisation results (60 trials, 5-fold
inner CV, AUC-PR objective, no SMOTE).}
\label{tab:optuna}
\centering
\begin{tabular}{lcccc}
\toprule
\textbf{Model} & \textbf{Best AUC-PR} & \textbf{Trials} &
\textbf{SMOTE} & \textbf{Wall Time (s)} \\
\midrule
CatBoost            & 0.7599 & 60 & No & 422  \\
GradientBoosting    & 0.7597 & 60 & No & 1137 \\
XGBoost             & 0.7572 & 60 & No & 236  \\
LightGBM            & 0.7480 &  1 & No & 3716 \\
RandomForest        & 0.7397 & 60 & No & 225  \\
LogisticRegression  & 0.7261 & 60 & No & 260  \\
\bottomrule
\end{tabular}
\end{table}

CatBoost achieved the highest inner-CV AUC-PR of 0.7599, marginally
ahead of GradientBoosting (0.7597) and XGBoost (0.7572). LightGBM's
Optuna run terminated early due to aggressive pruning, yielding only a
single completed trial; despite this, the pruned configuration still
attained 0.7480. LogisticRegression, as expected for a linear model on
this data, posted the lowest tuning score of 0.7261. The tuning
wall-times reveal an important practical trade-off: XGBoost (236 s) and
RandomForest (225 s) are substantially faster to optimise than
GradientBoosting (1137 s) or LightGBM (3716 s), making them preferable
in resource-constrained production environments.

% --------------------------------------------------------------------
\subsection{Model Analysis and Performance Evaluation}
% --------------------------------------------------------------------

Following tuning, the five tree-based models were carried forward for
full 10-fold stratified cross-validation (LogisticRegression was excluded
from downstream evaluation due to its substantially lower tuning score).
A weighted ensemble combining these five base models was also evaluated.
Table~\ref{tab:cv} reports the mean F1-score across all ten folds.

\begin{table}[H]
\caption{Ten-fold stratified cross-validation results: mean F1-score
for all model configurations (tuned, no SMOTE).}
\label{tab:cv}
\centering
\begin{tabular}{lcc}
\toprule
\textbf{Model} & \textbf{F1 (mean)} & \textbf{F1 (std)} \\
\midrule
WeightedEnsemble     & \textbf{0.7689} & 0.0334 \\
XGBoost              & 0.7350 & 0.0475 \\
LightGBM             & 0.7319 & 0.0274 \\
CatBoost             & 0.7249 & 0.0558 \\
GradientBoosting     & 0.7133 & 0.0554 \\
RandomForest         & 0.5860 & 0.0531 \\
\bottomrule
\end{tabular}
\end{table}

The weighted ensemble achieves the highest F1 of 0.7689, followed by
XGBoost (0.7350), LightGBM (0.7319), CatBoost (0.7249), and
GradientBoosting (0.7133). RandomForest lags substantially behind the
other models at F1~=~0.5860, suggesting that its bagging strategy is
less effective than boosting approaches on this dataset. Among the
individual models, XGBoost, LightGBM, and CatBoost form a competitive
tier separated by less than one percentage point of F1, while
GradientBoosting trails by roughly two points. The weighted ensemble
combines complementary strengths across these models to achieve the
best overall performance.
Figure~\ref{fig:model_comparison} provides a visual summary of the F1-scores across all models.

\begin{figure}[!htbp]
\centering
\includegraphics[width=0.85\textwidth]{fig_model_comparison.png}
\caption{Model comparison: F1-scores across all model configurations.}
\label{fig:model_comparison}
\end{figure}

\begin{figure}[htbp]
\centering
\includegraphics[width=0.6\textwidth]{fig_confusion_matrix.png}
\caption{Confusion matrix for the weighted ensemble at fixed threshold $t = 0.444$.}
\label{fig:confusion_matrix}
\end{figure}


% --------------------------------------------------------------------
\subsection{Statistical Significance}
\label{sec:stats}
% --------------------------------------------------------------------

To validate that observed performance differences are not attributable
to random variation across folds, a non-parametric Friedman test was
applied over the per-fold F1 scores of all six model configurations
(five base models plus the weighted ensemble).
The test returned $\chi^2 = 18.72$ with $p < 0.001$, rejecting
the null hypothesis that all models perform identically.

Post-hoc pairwise Wilcoxon signed-rank tests (with Holm-Bonferroni
correction) were then conducted with the weighted ensemble as the
reference model. The ensemble was found to be statistically
significantly better than:

\begin{itemize}[nosep]
  \item LightGBM ($p = 0.010$)
  \item GradientBoosting ($p = 0.010$)
  \item RandomForest ($p = 0.010$)
  \item XGBoost ($p = 0.010$)
  \item CatBoost ($p = 0.014$)
\end{itemize}

All five base models show statistically significant differences from
the weighted ensemble after Holm-Bonferroni correction, confirming
that the ensemble's improved F1 is not attributable to chance.
Figure~\ref{fig:stat_sig} visualises the adjusted p-values for all pairwise comparisons.

\begin{figure}[!htbp]
\centering
\includegraphics[width=0.85\textwidth]{fig_statistical_significance.png}
\caption{Pairwise Wilcoxon signed-rank test results with Holm-Bonferroni correction. Green bars indicate statistically significant differences ($p < 0.05$).}
\label{fig:stat_sig}
\end{figure}


% --------------------------------------------------------------------
\subsection{Threshold Optimisation}
% --------------------------------------------------------------------

The decision threshold was selected by optimising F1 on the validation
folds and then fixed at $t = 0.444$ for all downstream evaluation.
Table~\ref{tab:threshold} reports the weighted ensemble classification
metrics at this threshold.

\begin{table}[H]
\caption{Weighted ensemble classification metrics at fixed threshold $t = 0.444$.}
\label{tab:threshold}
\centering
\begin{tabular}{lccc}
\toprule
\textbf{Threshold} & \textbf{F1} & \textbf{Precision} & \textbf{Recall} \\
\midrule
0.444 (fixed)   & 0.7345 & 0.7368 & 0.7470 \\
\bottomrule
\end{tabular}
\end{table}

Fixing the threshold at the validation-optimal value of 0.444, rather
than the default 0.5, yields a slight recall improvement while
maintaining precision. This threshold was adopted for all downstream
bootstrap and calibration analyses.

\begin{figure}[!htbp]
  \centering
  \includegraphics[width=0.80\textwidth]{fig_threshold_sensitivity.png}
  \caption{Threshold sensitivity curve for the weighted ensemble: F1,
  Precision, and Recall as a function of decision threshold. The
  vertical dashed line marks the fixed threshold at $t = 0.444$.}
  \label{fig:threshold}
\end{figure}

% --------------------------------------------------------------------
\subsection{Bootstrap Confidence Intervals}
% --------------------------------------------------------------------

To quantify estimation uncertainty, a stratified bootstrap with
$n = 2{,}000$ resamples was applied to the weighted ensemble's test-set
predictions at the fixed threshold $t = 0.444$. The resulting 95\%
confidence intervals are reported in Table~\ref{tab:bootstrap}.

\begin{table}[H]
\caption{Bootstrap confidence intervals for the weighted ensemble at
fixed threshold $t = 0.444$ ($n = 2{,}000$ resamples, 95\% CI).}
\label{tab:bootstrap}
\centering
\begin{tabular}{lccc}
\toprule
\textbf{Metric} & \textbf{Point Estimate} &
\textbf{95\% CI Lower} & \textbf{95\% CI Upper} \\
\midrule
F1         & 0.7416 & 0.7118 & 0.7715 \\
Precision  & 0.7369 & 0.6985 & 0.7750 \\
Recall     & 0.7467 & 0.7094 & 0.7839 \\
MCC        & 0.7148 & 0.6818 & 0.7474 \\
ROC-AUC    & 0.9586 & 0.9477 & 0.9686 \\
PR-AUC     & 0.8115 & 0.7801 & 0.8406 \\
\bottomrule
\end{tabular}
\end{table}

All confidence intervals are reasonably narrow, reflecting stable
generalisation. The ROC-AUC of $0.9586\;[0.9477,\,0.9686]$ confirms
strong discriminative power across the full threshold range, while
the PR-AUC of $0.8115\;[0.7801,\,0.8406]$ is particularly informative
given the class imbalance: a random classifier would score near the
minority-class prevalence ($\approx 0.094$), so 0.8115 represents a
substantial lift of more than eight times the baseline. The bootstrapped
MCC of $0.7148\;[0.6818,\,0.7474]$ corroborates the F1-based findings
without being susceptible to the known sensitivity of F1 to class
imbalance.

\begin{figure}[htbp]
\centering
\includegraphics[width=0.75\textwidth]{fig_bootstrap_ci.png}
\caption{Bootstrap 95\% confidence intervals for the weighted ensemble (2,000 iterations, fixed threshold $t=0.444$).}
\label{fig:bootstrap_ci}
\end{figure}

\begin{figure}[htbp]
\centering
\includegraphics[width=0.65\textwidth]{fig_pr_curve.png}
\caption{Precision-recall curve for the weighted ensemble (PR-AUC = 0.8114). The operating point at fixed threshold $t = 0.444$ achieves precision = 0.737 and recall = 0.747. The baseline precision (dashed) corresponds to the fraud prevalence of 9.4\%.}
\label{fig:pr_curve}
\end{figure}

% --------------------------------------------------------------------
\subsection{SMOTE-ENN Ablation Study}
% --------------------------------------------------------------------

To isolate the effect of synthetic oversampling, XGBoost, LightGBM, and
GradientBoosting were each re-evaluated under two conditions: with
SMOTE-ENN applied inside each training fold, and without. Table~\ref{tab:smote}
summarises the F1 outcome.

\begin{table}[H]
\caption{SMOTE-ENN ablation: F1 scores with and without oversampling
(10-fold CV, tuned models).}
\label{tab:smote}
\centering
\begin{tabular}{lccc}
\toprule
\textbf{Model} & \textbf{With SMOTE-ENN} & \textbf{No SMOTE} &
\textbf{Change} \\
\midrule
XGBoost          & 0.6639 & \textbf{0.7350} & $-9.7\%$ \\
LightGBM         & 0.6815 & \textbf{0.7319} & $-6.9\%$  \\
GradientBoosting & 0.6528 & \textbf{0.7133} & $-8.5\%$ \\
\bottomrule
\end{tabular}
\end{table}

Across all three tree-based ensembles, SMOTE-ENN oversampling
consistently \emph{degrades} F1 - by 9.7\% for XGBoost, 8.5\% for
GradientBoosting, and 6.9\% for LightGBM. This counter-intuitive result
is well documented in the literature for gradient-boosted trees: the
native loss-weighted splitting in XGBoost (via \texttt{scale\_pos\_weight})
and the analogue mechanisms in LightGBM and GradientBoosting already
account for class imbalance. Introducing SMOTE-ENN on top of these
mechanisms injects synthetic minority-class samples into the training
distribution, distorting the decision boundary and inflating recall
while degrading precision. The net effect is a lower F1 and, more
critically, an inflated false-positive rate that would be operationally
costly in a claims-audit workflow. Figure~\ref{fig:smote_abl} illustrates this finding. Consequently, no oversampling is used
in the final deployed pipeline.

\begin{figure}[!htbp]
\centering
\includegraphics[width=0.82\textwidth]{fig_smote_ablation.png}
\caption{SMOTE-ENN ablation: F1-score comparison with and without synthetic oversampling for three gradient-boosted models. Delta percentages indicate the performance penalty introduced by SMOTE-ENN.}
\label{fig:smote_abl}
\end{figure}

% --------------------------------------------------------------------
\subsection{Feature Importance Ranking}
% --------------------------------------------------------------------

Feature importance was assessed through two complementary approaches:
SHAP (SHapley Additive exPlanations) values, which provide
theoretically grounded attribution based on cooperative game theory,
and permutation importance, which measures the drop in performance
when each feature's values are randomly shuffled.

Figure~\ref{fig:shap_beeswarm} presents the SHAP beeswarm plot for the
tuned XGBoost model, displaying the distribution of SHAP values for the
top 20 features across all test instances. Each point represents one
provider, coloured by feature value (red = high, blue = low).

\begin{figure}[!htbp]
  \centering
  \includegraphics[width=0.90\textwidth]{fig_shap_beeswarm_tuned.png}
  \caption{SHAP beeswarm plot (XGBoost): distribution of SHAP
  values for the top 20 features across all test-set providers. Red
  points indicate high feature values; blue points indicate low values.}
  \label{fig:shap_beeswarm}
\end{figure}

Figure~\ref{fig:shap_bar} shows the mean absolute SHAP values (global
importance ranking) in a bar chart, providing a concise summary of
which features drive model predictions most strongly.

\begin{figure}[!htbp]
  \centering
  \includegraphics[width=0.85\textwidth]{fig_shap_bar_tuned.png}
  \caption{SHAP global feature importance (XGBoost): mean
  absolute SHAP values ranked in descending order.}
  \label{fig:shap_bar}
\end{figure}

Figure~\ref{fig:shap_dep} presents SHAP dependence plots for the four
most important features, illustrating how the model's output changes
as each feature varies, with interaction effects encoded by the
secondary colour axis.

\begin{figure}[!htbp]
  \centering
  \includegraphics[width=0.95\textwidth]{fig_shap_dependence_tuned.png}
  \caption{SHAP dependence plots for the top-4 features (XGBoost).
  Colour represents the value of the strongest interaction
  feature identified by SHAP.}
  \label{fig:shap_dep}
\end{figure}

Figure~\ref{fig:perm_imp} shows permutation importance, confirming the
SHAP-derived ranking through a model-agnostic lens and providing
error bars over 10 shuffling repetitions.

\begin{figure}[!htbp]
  \centering
  \includegraphics[width=0.85\textwidth]{fig_permutation_importance.png}
  \caption{Permutation feature importance (XGBoost): mean
  decrease in F1 when each feature is randomly permuted, with
  error bars over 10 repetitions.}
  \label{fig:perm_imp}
\end{figure}

% --------------------------------------------------------------------
\subsection{SHAP Waterfall Case Studies}
% --------------------------------------------------------------------

To illustrate how the model reasons about individual providers, SHAP
waterfall plots were generated for three clinically meaningful cases: a
true positive (TP), a false positive (FP), and a false negative (FN).
Each waterfall chart decomposes the gap between the model's base value
and the final prediction into additive feature contributions.

Figure~\ref{fig:wf_tp} shows the true-positive case, where the model
correctly identifies a fraudulent provider. The dominant positive
contributions align with the globally important features identified in
Section 4.2.8, confirming internal consistency.

\begin{figure}[!htbp]
  \centering
  \includegraphics[width=0.85\textwidth]{fig_shap_waterfall_tp_tuned.png}
  \caption{SHAP waterfall plot - True Positive case (XGBoost).
  The model correctly predicts fraud; positive SHAP values push the
  prediction above the decision threshold.}
  \label{fig:wf_tp}
\end{figure}

Figure~\ref{fig:wf_fp} examines the false-positive case, where the
model incorrectly flags a legitimate provider. Understanding which
features mislead the model in these cases is important for refining
the feature-engineering pipeline and for communicating audit priorities
to investigators.

\begin{figure}[!htbp]
  \centering
  \includegraphics[width=0.85\textwidth]{fig_shap_waterfall_fp_tuned.png}
  \caption{SHAP waterfall plot - False Positive case (XGBoost).
  The model incorrectly predicts fraud for a legitimate provider,
  revealing which features are responsible for the erroneous signal.}
  \label{fig:wf_fp}
\end{figure}

Figure~\ref{fig:wf_fn} shows the false-negative case, where the model
fails to flag an actual fraudulent provider. Analysing the suppressive
SHAP contributions in this case highlights potential weaknesses in the
current feature set or thresholding strategy.

\begin{figure}[!htbp]
  \centering
  \includegraphics[width=0.85\textwidth]{fig_shap_waterfall_fn_tuned.png}
  \caption{SHAP waterfall plot - False Negative case (XGBoost).
  The model misses an actual fraudulent provider; negative SHAP
  contributions suppress the prediction below the decision threshold.}
  \label{fig:wf_fn}
\end{figure}

% --------------------------------------------------------------------
\subsection{LIME Explanations}
% --------------------------------------------------------------------

LIME (Local Interpretable Model-agnostic Explanations) was applied as
an independent local-explanation technique to cross-validate the SHAP
waterfall findings. LIME fits a weighted linear model in the neighbourhood
of each instance, providing a complementary, model-agnostic perspective
on individual predictions.

Figure~\ref{fig:lime_tp} shows the LIME explanation for the same
true-positive provider examined in the SHAP waterfall analysis.
Features coloured in orange support the fraud prediction; features in
blue oppose it.

\begin{figure}[!htbp]
  \centering
  \includegraphics[width=0.85\textwidth]{fig_lime_tp_tuned.png}
  \caption{LIME local explanation - True Positive case (XGBoost).
  Orange bars indicate features supporting the fraud prediction;
  blue bars indicate features opposing it.}
  \label{fig:lime_tp}
\end{figure}

Figure~\ref{fig:lime_fp} presents the LIME explanation for the
false-positive provider. Comparing these weights against the
corresponding SHAP waterfall (Figure~\ref{fig:wf_fp}) reveals whether
both explanation methods agree on the misleading features.

\begin{figure}[!htbp]
  \centering
  \includegraphics[width=0.85\textwidth]{fig_lime_fp_tuned.png}
  \caption{LIME local explanation - False Positive case (XGBoost).
  The linear approximation reveals which feature values drove the
  erroneous fraud classification.}
  \label{fig:lime_fp}
\end{figure}

Figure~\ref{fig:lime_fn} shows the LIME explanation for the
false-negative provider. Where LIME and SHAP agree on the suppressive
features, this provides stronger evidence that the model has learned a
systematic blind spot for this provider sub-type.

\begin{figure}[!htbp]
  \centering
  \includegraphics[width=0.85\textwidth]{fig_lime_fn_tuned.png}
  \caption{LIME local explanation - False Negative case (XGBoost).
  Features with negative LIME weights suppress the fraud score,
  explaining why the model misses this fraudulent provider.}
  \label{fig:lime_fn}
\end{figure}

% --------------------------------------------------------------------
\subsection{Feature Selection Ablation}
% --------------------------------------------------------------------

To determine the minimum feature subset that preserves predictive
performance, a stepwise feature-selection ablation was conducted.
Starting from the top-10 SHAP-ranked features and incrementally adding
features up to the full 190-feature set, XGBoost was re-evaluated via
10-fold CV at each step. Results are summarised in Table~\ref{tab:featsel}.

\begin{table}[H]
\caption{Feature-selection ablation: XGBoost F1 as a function of
top-$K$ SHAP-ranked features (10-fold CV, tuned).}
\label{tab:featsel}
\centering
\begin{tabular}{cc}
\toprule
\textbf{Feature Subset Size ($K$)} & \textbf{F1} \\
\midrule
10  & 0.6534 \\
20  & 0.6763 \\
30  & 0.6716 \\
40  & 0.6711 \\
190 & \textbf{0.7350} \\
\bottomrule
\end{tabular}
\end{table}

A substantial performance gap separates the reduced subsets from the
full feature set. Using only the top 10 features yields F1 = 0.6534,
a gap of 8.2 percentage points below the full-feature F1 of 0.7350.
Increasing the subset to 20 features recovers some performance
(F1 = 0.6763), but further expansion to 30 and 40 features does not
continue this improvement monotonically - suggesting diminishing returns
and possible feature redundancy in the 21-40 rank range. The full
feature set consistently outperforms all reduced variants, indicating
that the less-prominent features (ranks 11-190) collectively encode
complementary signals that the model exploits. This result argues
against aggressive feature pruning in the current pipeline; future
work could explore structured feature compression (e.g., PCA or
autoencoders) rather than SHAP-rank truncation. Figure~\ref{fig:feat_sel} visualises this trend.

\begin{figure}[!htbp]
\centering
\includegraphics[width=0.75\textwidth]{fig_feature_selection.png}
\caption{Feature selection ablation: F1-score as a function of the number of top-K features selected by permutation importance. The full 190-feature set achieves the best performance.}
\label{fig:feat_sel}
\end{figure}

% --------------------------------------------------------------------
\subsection{Learning Curves}
% --------------------------------------------------------------------

\begin{figure}[!htbp]
  \centering
  \includegraphics[width=0.85\textwidth]{fig_learning_curves.png}
  \caption{Learning curves for XGBoost: training F1 and
  cross-validation F1 as a function of training-set size. The
  plateau in validation F1 and the gap between training and
  validation curves indicate the degree of overfitting.}
  \label{fig:learning}
\end{figure}

The learning curves in Figure~\ref{fig:learning} reveal a classic
overfitting signature: training F1 reaches 1.0 across all training-set
sizes, while validation F1 plateaus near 0.73. The large and persistent
gap between the two curves confirms that XGBoost has sufficient capacity
to perfectly memorise the training data but that the learned decision
boundary does not generalise fully to held-out providers. This behaviour
is expected for gradient-boosted trees with deep individual estimators
on a moderately sized tabular dataset. The validation curve
does not exhibit a downward trend as training size increases, indicating
that the model is not in a high-variance regime where more data would
hurt. Additional training data (e.g., incorporating Medicare Part A or
DMEPOS claims) could plausibly narrow this gap. Regularisation
parameters such as \texttt{min\_child\_weight}, \texttt{max\_depth},
and \texttt{subsample} were already tuned by Optuna, so the remaining
gap reflects an irreducible component related to feature informativeness
rather than model over-complexity.

% --------------------------------------------------------------------
\subsection{Calibration Analysis}
% --------------------------------------------------------------------

\begin{figure}[!htbp]
  \centering
  \includegraphics[width=0.80\textwidth]{fig_calibration.png}
  \caption{Probability calibration curve for XGBoost (reliability
  diagram): mean predicted probability versus fraction of positives in
  each probability bin. A perfectly calibrated model lies on the
  diagonal. The histogram shows the distribution of predicted
  probabilities.}
  \label{fig:calibration}
\end{figure}

Figure~\ref{fig:calibration} plots the reliability diagram for XGBoost.
The model is well-calibrated in the low-probability range (0-0.3)
but tends to over-predict fraud probability in the mid-range (0.3-0.7),
a pattern common to gradient-boosted trees that produce overconfident
probability estimates near the decision boundary. In the high-confidence
region ($p > 0.7$), calibration recovers. For operational deployment,
Platt scaling or isotonic regression post-processing could be applied
to reduce mid-range over-confidence and improve the reliability of the
fraud probability scores displayed on the dashboard, particularly when
risk-level boundaries (\textit{Low}/\textit{Medium}/\textit{High}) are
defined by fixed probability cutoffs.


% --------------------------------------------------------------------
\subsection{Blockchain Verification}
\label{sec:blockchain_verification}
% --------------------------------------------------------------------

Upon completion of the full pipeline run, all blockchain integrity
checks were executed programmatically. The results are as follows:

\begin{itemize}[nosep,leftmargin=1.5em]
  \item \textbf{Chain length:} 110 blocks, 5{,}410 provider records
        registered.
  \item \textbf{Chain validation:} PASSED - all sequential
        \texttt{previous\_hash} links are intact and every block hash
        satisfies the SHA-256 Proof-of-Work constraint (difficulty 2,
        leading zeros verified).
  \item \textbf{Merkle proof:} PASSED - membership proof traversed
        6 Merkle-tree levels successfully; root digest matches the
        stored \texttt{merkle\_root} field in the block header.
  \item \textbf{ECIES decryption:} PASSED - all \texttt{provider\_id}
        fields were successfully decrypted using the private key,
        confirming round-trip encryption integrity.
  \item \textbf{Block build time:} 7.0 seconds for 5{,}410 records,
        averaging approximately 1.3 ms per record - well within
        real-time throughput requirements for claims-audit workflows.
  \item \textbf{Tamper detection:} A synthetic tamper test (modifying
        a single field in block 50) was detected immediately at the
        next hash-chain validation step, confirming auditability.
\end{itemize}

These results confirm that the blockchain layer provides end-to-end
data integrity, privacy-preserving storage, and cryptographic
auditability without introducing unacceptable latency. The combination
of SHA-256 chaining, Merkle proofs, and ECIES encryption satisfies
the transparency and non-repudiation requirements stipulated for
healthcare fraud audit trails under relevant regulatory frameworks.
